\Needspace{10\baselineskip}
\centerline{\textsc{Ceres.}}

\medskip
Monatliche Correspondenz. Band VIII, S.\ 288—291 und S.\ 369—371. September und Oktober 1803.

\medskip
— — —

\medskip
Dr.\ \textsc{Gauß} untersuchte unterdessen, welche Korrektur seiner letzten Elemente (VIII) sich aus den Bremer Beobachtungen und den drei auf S.\ 94 des Julihefts angegebenen Beobachtungen aus Palermo ergab. Eine kleine Änderung genügte, um diese neuen Beobachtungen mit den alten in Einklang zu bringen. Nur den Knoten musste er um 3 Minuten weiter verschieben als für die Beobachtungen von 1801 in Palermo. Hoffentlich lässt sich dies später durch eine vollständigere Berechnung der Störungen rechtfertigen. Hier sind nun diese neuen Elemente (IX).

\begin{center}
\begingroup\fontsize{9}{11}\selectfont\setlength{\tabcolsep}{3pt}
\begin{tabular}{lr}
Epoche für Seeberg 1803 . . . . . . . . . . . . & $233^\circ\ 36'\ 3''1$ \\
Tropische Bewegung in 365 Tagen . . . . . . . . . & $78\ 8\ 7.2$ \\
Tropische Bewegung an einem Tag . . . . . . . . & $770''650$ \\
Logarithmus der großen Halbachse . . . . . . . . & $0.4421516$ \\
Aphel 1803 . . . . . . . . . . . . . . . & $326^\circ\ 33'\ 18''$ \\
Knoten 1803 . . . . . . . . . . . . . . . . . & $80\ 58\ 22$ \\
Exzentrizität . . . . . . . . . . . . . . . . & $0.0788941$ \\
Neigung . . . . . . . . . . . . . . . . . . & $10^\circ\ 37'\ 54''$ \\
\end{tabular}\endgroup
\end{center}
\SourceLeaf{251}%

Diese Elemente stimmen mit den drei erwähnten Beobachtungen von \textsc{Piazzi} wie folgt überein:

\begin{center}
\begingroup\fontsize{9}{11}\selectfont\setlength{\tabcolsep}{3pt}\renewcommand{\arraystretch}{1.12}
\begin{tabular}{r|r|r|r|r|r}
 & & \multicolumn{2}{c|}{Berechnet} & \multicolumn{2}{c}{Unterschied R—B} \\ \cline{3-6}
1803 & \shortstack{Mittlere Zeit\\in Palermo} & \shortstack{Rektaszension\\der $\CeresSymbol$} & \shortstack{Deklination\\der $\CeresSymbol$} & in Rekt. & in Dekl. \\ \hline
Mai 12 & $15^u\,53^m\,37^s7$ & $288^\circ\,19'\,17''9$ & $24^\circ\,36'\,22''0$ & $+\,2''9$ & $+\,9''0$ \\
--- 13 & $15\,49\,41.3$ & $288\,18\,57.3$ & $24\,39\,57.2$ & $-\,10.2$ & $-\,4.2$ \\
--- 14 & $15\,45\,41.9$ & $288\,18\,13.1$ & $24\,43\,36.7$ & $-\,1.9$ & $+\,5.6$ \\
\end{tabular}\endgroup
\end{center}

[Das Oktoberheft enthält die nach den Elementen IX von \textsc{Gauß} berechnete Ephemeride. Sie reicht vom 30. April 1804 bis zum 19. Januar 1805 und gilt in Abständen von drei Tagen jeweils für Mitternacht in Seeberg. Sie gibt die Rektaszension und die Deklination auf Minuten genau, die Lichtstärke auf fünf Dezimalstellen und auch die mittlere Zeit des Meridiandurchgangs an.]

\begin{flushright}
[\textsc{von Zach}.]
\end{flushright}

\bigskip
